\section{Proofs}
\label{app:proofs}

\subsection{Finite-set JL projection}

\begin{proof}
For a fixed nonzero vector \(v\), concentration of
\(\norm{\Pi v}_2^2/\norm v_2^2\) around one gives
\[
    \Prob\!\left(
      \abs{\norm{\Pi v}_2^2-\norm v_2^2}
      >\varepsilon\norm v_2^2
    \right)
    \leq 2e^{-c_0\varepsilon^2n}
\]
for a universal \(c_0>0\). Apply this inequality to the fewer than \(M^2/2\)
nonzero differences \(u-v\) and take a union bound. Choosing
\(n\geq C\varepsilon^{-2}\log(M/\delta)\) for a sufficiently large universal
\(C\) makes the total failure probability at most \(\delta\). This is the
standard finite-set JL argument \citep{DasguptaGupta2003,Woodruff2014}. If
\(\mathcal P\) is random but independent of \(\Pi\), the calculation holds
conditionally on every realized \(\mathcal P\), and hence also unconditionally.
\end{proof}

\subsection{Trace-norm contraction}

\begin{proof}
For density operators \(\Tr[\rho]=\Tr[\rho']=1\), so the replacement
components cancel when two outputs with the same input are subtracted:
\[
    T_{r,x}(\rho)-T_{r,x}(\rho')
    =
    \lambda_rV_r(x)(\rho-\rho')V_r(x)^\dagger.
\]
Homogeneity and unitary invariance of the trace norm give
\cref{eq:exact-contraction}.

Iterating the equality for \(w\) common inputs yields
\[
    \traceNorm{\rho_0^{(r)}-\rho_0^{\prime(r)}}
    \leq
    \lambda_r^w
    \traceNorm{\rho_{-w}^{(r)}-\rho_{-w}^{\prime(r)}}
    \leq2\lambda_r^w.
\]
To construct the causal state, initialize the recursion at time \(-m\) in an
arbitrary density operator and call the resulting time-zero state
\(\rho_0^{[m]}\). If \(m'>m\), regard the recursion started at \(-m'\) as
having some density operator at time \(-m\). The last \(m\) updates are common,
so
\[
    \traceNorm{\rho_0^{[m']}-\rho_0^{[m]}}
    \leq2\lambda_r^m.
\]
Thus \((\rho_0^{[m]})_m\) is Cauchy in the finite-dimensional trace-norm
space and converges. The bound also shows that the limit is independent of
the chosen initial states. Shifting the argument in time produces a unique
causal state sequence. For every finite start time, the time-zero state is a
measurable function of a finite input block. Under
\cref{sec:analytical-scope}, the causal state is their pointwise trace-norm
limit and is therefore measurable.
\end{proof}

\subsection{Finite-window feature error}

\begin{proof}
For every \(r\) and \(O\in\cO\), \cref{lem:contraction} and trace duality give
\[
    \abs{
      \Tr\!\left[
        O\bigl(\rho_{\infty,0}^{(r)}-\rho_{w,0}^{(r)}\bigr)
      \right]
    }
    \leq
    \opNorm{O}
    \traceNorm{\rho_{\infty,0}^{(r)}-\rho_{w,0}^{(r)}}
    \leq2\lambda_r^w.
\]
Squaring and summing over \(r\) and \(O\) yields
\[
    \norm{\PhiInf-\PhiWin}_2^2
    \leq4|\cO|\sum_{r=1}^R\lambda_r^{2w}.
\]
Taking square roots proves the first inequality. The second follows from
\(\lambda_r\leq\lambdaMax\).
\end{proof}

\subsection{Mat\'ern RKHS section stability}

\begin{proof}
The reproducing property and normalization give
\begin{align*}
    \norm{\kappa(z,\cdot)-\kappa(z',\cdot)}_{\cH_\kappa}^2
    &=
    \kappa(z,z)+\kappa(z',z')-2\kappa(z,z')\\
    &=
    2\left(1-\varphi_{\nu,\xi}(\norm{z-z'}_2)\right).
\end{align*}
Let \(u=z-z'\). Using the spectral representation in
\cref{eq:matern-bochner}, the elementary inequality
\(1-\cos t\leq t^2/2\), and
\cref{eq:matern-spectral-covariance},
\begin{align*}
    2(1-\kappa(z,z'))
    &=
    2\E[1-\cos(\Omega^\top u)]\\
    &\leq
    \E[(\Omega^\top u)^2]\\
    &=
    u^\top\E[\Omega\Omega^\top]u\\
    &=
    \frac{\nu}{(\nu-1)\xi^2}\norm u_2^2.
\end{align*}
Taking square roots proves the claim.
\end{proof}

\subsection{Prediction and loss truncation}

\begin{proof}
For any \(h\in\Hball\), the reproducing property,
\cref{lem:matern-stability}, and \cref{prop:feature-truncation} imply
\begin{align*}
    \abs{h(\PhiInf)-h(\PhiWin)}
    &=
    \abs{
      \ip{h}{
        \kappa(\PhiInf,\cdot)-\kappa(\PhiWin,\cdot)
      }_{\cH_\kappa}
    }\\
    &\leq
    \Lambda
    \norm{
      \kappa(\PhiInf,\cdot)-\kappa(\PhiWin,\cdot)
    }_{\cH_\kappa}\\
    &\leq
    \frac{2\Lambda}{\xi}
    \sqrt{
      \frac{\nu|\cO|}{\nu-1}
      \sum_{r=1}^R\lambda_r^{2w}
    }.
\end{align*}
This is \cref{eq:prediction-truncation}.

For \(u,v\in[-\Lambda,\Lambda]\) and \(\abs y\leq\Upsilon\),
\[
    \abs{(u-y)^2-(v-y)^2}
    =
    \abs{u-v}\abs{u+v-2y}
    \leq2(\Lambda+\Upsilon)\abs{u-v}.
\]
Both predictions are in \([-\Lambda,\Lambda]\) by
\cref{eq:prediction-bound}. Apply the last display pointwise with
\(u=h(\PhiInf)\), \(v=h(\PhiWin)\), use the prediction bound above, and take
expectations. This gives
\(\abs{\risk_\infty(h)-\risk_w(h)}\leq\tau_w\).
\end{proof}

\subsection{Conditional interpolation}

\begin{proof}
The Mat\'ern spectral density in \cref{eq:matern-spectral-density} is strictly
positive on \(\R^p\). For any \(\bm c\in\R^N\), Bochner's representation gives
\[
    \bm c^\top K\bm c
    =
    \int_{\R^p}
    \left|\sum_{i=1}^N c_i e^{\ii\omega^\top z_i}\right|^2
    p_{\nu,\xi}(\omega)\,d\omega.
\]
Suppose that the \(z_i\) are distinct and \(\bm c\neq0\). Choose a direction
\(v\in\R^p\) outside the finite union of hyperplanes
\(\{v:v^\top(z_i-z_j)=0\}\), so the scalars \(v^\top z_i\) are pairwise
distinct. The restriction
\(t\mapsto\sum_i c_i e^{\ii t v^\top z_i}\) is a one-dimensional exponential
polynomial with distinct frequencies and cannot vanish identically unless all
\(c_i=0\). Thus the multivariate exponential polynomial is nonzero at some
point and, by continuity, on an open set of positive Lebesgue measure. Since
\(p_{\nu,\xi}>0\) everywhere, the integral is strictly positive. Hence \(K\)
is positive definite and invertible.

Let \(\bm\alpha=K^{-1}\bm y\) and
\[
    h_\star(\cdot)=\sum_{i=1}^N\alpha_i\kappa(z_i,\cdot).
\]
Then \((h_\star(z_1),\ldots,h_\star(z_N))^\top=K\bm\alpha=\bm y\) and
\[
    \norm{h_\star}_{\cH_\kappa}^2
    =\bm\alpha^\top K\bm\alpha
    =\bm y^\top K^{-1}\bm y
    =\Lambda_\star^2.
\]
Moreover, if \(h\) is any other interpolant, then
\(g:=h-h_\star\) satisfies \(g(z_i)=0\) for all \(i\), so
\(g\) is orthogonal to every kernel section \(\kappa(z_i,\cdot)\) and hence to
\(h_\star\). Therefore
\(\norm{h}_{\cH_\kappa}^2=\norm{h_\star}_{\cH_\kappa}^2+
\norm{g}_{\cH_\kappa}^2\), proving that \(h_\star\) is the unique minimum-norm
interpolant \citep{Aronszajn1950,KimeldorfWahba1971}.

Thus \(h_\star\) is feasible for every \(\Lambda\geq\Lambda_\star\) and has
zero empirical risk. For \(0\leq\Lambda\leq\Lambda_\star\), the function
\(h_\Lambda=(\Lambda/\Lambda_\star)h_\star\) is feasible and satisfies
\[
    h_\Lambda(z_i)-y_i
    =-\left(1-\frac{\Lambda}{\Lambda_\star}\right)y_i.
\]
Its empirical risk is the right-hand side of
\cref{eq:conditional-interpolation}; the minimum over the ball cannot be
larger. The case \(\bm y=0\) is immediate.
\end{proof}

\subsection{Mixing of the window process}

\begin{proof}
By stationarity it is enough to compare the sigma-algebra generated by
\((W_j)_{j\leq0}\) with that generated by \((W_j)_{j\geq a}\). Every raw
variable occurring in a past window has time index at most zero, so
\[
    \sigma(W_j:j\leq0)
    \subseteq
    \sigma(Z_t:t\leq0).
\]
The first raw input appearing in \(W_a\) has index \(as-w+1\), and all raw
variables in later windows occur no earlier. Therefore
\[
    \sigma(W_j:j\geq a)
    \subseteq
    \sigma(Z_t:t\geq as-w+1).
\]
Monotonicity of \(\beta(\cdot,\cdot)\) under sigma-algebra inclusion gives
\[
    \beta_W(a)
    \leq
    \beta_Z(as-w+1)
    \leq
    \beta_Z(as-w),
\]
where the last inequality uses that mixing coefficients are nonincreasing.
For \(a=1\) and \(s=w+g\), \(as-w=g\).
\end{proof}

\subsection{One-sided dependent-data risk bound}

\begin{proof}
For \(h\in\Hball\), define
\[
    f_h(W_j)
    =
    \left(
      h(\PhiWin(X_{js-w+1:js}))-Y_{js}
    \right)^2.
\]
By \cref{eq:prediction-bound} and \(\abs{Y_t}\leq\Upsilon\),
\[
    0\leq f_h(W_j)\leq B,
    \qquad B=(\Lambda+\Upsilon)^2.
\]
The class \(\cF_\Lambda=\set{f_h:h\in\Hball}\) is fixed independently of the
sample under the analytical scope of \cref{sec:analytical-scope}. Apply \cref{eq:mr-background} to the
stationary window process. By \cref{lem:window-mixing},
\[
    \beta_W(a)
    \leq\beta_Z(as-w),
\]
Let \(\delta'_W=\delta-4(\mu-1)\beta_W(a)\) denote the source-theorem
remainder. The mixing bound gives \(0<\delta'\leq\delta'_W\), so the source
theorem applies; using the smaller \(\delta'\) only enlarges the confidence
term. Consequently, simultaneously for all \(h\in\Hball\),
\[
    \risk_w(h)
    \leq
    \erisk_{N,w}(h)
    +\hRad_{\cS_\mu}(\cF_\Lambda)
    +3(\Lambda+\Upsilon)^2
      \sqrt{\frac{\log(4/\delta')}{2\mu}}.
\]
Finally, \cref{cor:loss-truncation} gives
\(\risk_\infty(h)\leq\risk_w(h)+\tau_w\). Combining the two inequalities
proves \cref{eq:empirical-rad-bound}.
\end{proof}

\subsection{Closed-form RKHS complexity bound}

\begin{proof}
Condition on \(\cS_\mu=((u_j,y_j))_{j=1}^\mu\), where
\(u_j=\PhiWin(x_j)\), and let \(\mathcal G_\Lambda\) be the prediction class
induced by \(\Hball\) on these features. Because the RKHS ball is symmetric,
\begin{align*}
    \hRad_{\cS_\mu}(\mathcal G_\Lambda)
    &=
    \E_{\bm\sigma}\!\left[
      \sup_{\norm{h}_{\cH_\kappa}\leq\Lambda}
      \abs{\frac2\mu
      \ip{h}{\sum_{j=1}^\mu\sigma_j\kappa(u_j,\cdot)}_{\cH_\kappa}}
    \right]\\
    &=
    \frac{2\Lambda}{\mu}
    \E_{\bm\sigma}
    \norm{\sum_{j=1}^\mu\sigma_j\kappa(u_j,\cdot)}_{\cH_\kappa}\\
    &\leq
    \frac{2\Lambda}{\mu}
    \sqrt{\sum_{j=1}^\mu\kappa(u_j,u_j)}
    =\frac{2\Lambda}{\sqrt\mu}.
\end{align*}

For each \(j\), define
\(\psi_j(t)=(t-y_j)^2-y_j^2\). On \([-\Lambda,\Lambda]\),
\(\psi_j(0)=0\) and \(\psi_j\) is
\(2(\Lambda+\Upsilon)\)-Lipschitz. The absolute-value contraction principle
\citep{BartlettMendelson2002} gives the safe bound
\[
    \E_{\bm\sigma}\sup_{h\in\Hball}
    \abs{\frac2\mu\sum_{j=1}^\mu\sigma_j\psi_j(h(u_j))}
    \leq
    4(\Lambda+\Upsilon)
    \hRad_{\cS_\mu}(\mathcal G_\Lambda)
    \leq
    \frac{8\Lambda(\Lambda+\Upsilon)}{\sqrt\mu}.
\]
The source convention contains an absolute value, so the label-only offset must
also be retained. By Jensen's inequality and independence of the Rademacher
signs,
\[
    \E_{\bm\sigma}
    \abs{\frac2\mu\sum_{j=1}^\mu\sigma_j y_j^2}
    \leq
    \frac2\mu\sqrt{\sum_{j=1}^\mu y_j^4}
    \leq\frac{2\Upsilon^2}{\sqrt\mu}.
\]
The triangle inequality therefore yields
\[
    \hRad_{\cS_\mu}(\cF_\Lambda)
    \leq
    \frac{8\Lambda(\Lambda+\Upsilon)+2\Upsilon^2}{\sqrt\mu}
    =\frac{C_{\Lambda,\Upsilon}}{\sqrt\mu}.
\]
Substitution into \cref{eq:empirical-rad-bound}, followed by the definition of
\(\tau_w\), proves \cref{eq:closed-generalization}.
\end{proof}

\subsection{Finite-family model selection}

\begin{proof}
For candidate \(\ell\), let
\[
    \cF_\ell
    :=
    \set{(x_{-w+1:0},y)\mapsto
      (h(\Phi_w^{(\ell)}(x_{-w+1:0}))-y)^2:
      h\in\cH_\ell},
\]
and define the fixed union
\(\cF_{\mathfrak C}:=\bigcup_{\ell=1}^L\cF_\ell\). Every function in this
union is measurable and takes values in \([0,B_{\mathfrak C}]\). Apply
\cref{eq:mr-background} once to \(\cF_{\mathfrak C}\), using
\cref{lem:window-mixing} exactly as in the proof of
\cref{thm:generalization}. With probability at least \(1-\delta\),
simultaneously for every \(\ell\) and \(h\in\cH_\ell\),
\begin{equation}
    \risk_{w,\ell}(h)
    \leq
    \erisk_{N,w,\ell}(h)
    +\hRad_{\cS_\mu}(\cF_{\mathfrak C})
    +3B_{\mathfrak C}
      \sqrt{\frac{\log(4/\delta')}{2\mu}}.
    \label{eq:finite-union-intermediate}
\end{equation}

Condition on \(\cS_\mu\) and define
\[
    Z_\ell(\bm\sigma)
    :=
    \sup_{f\in\cF_\ell}
    \abs{\frac2\mu\sum_{j=1}^\mu
    \sigma_j f(\widetilde W_j)}.
\]
Changing one Rademacher sign changes \(Z_\ell\) by at most
\(4B_{\mathfrak C}/\mu\). The bounded-difference exponential inequality
\citep{McDiarmid1989} therefore gives, for every \(t>0\),
\[
    \E_{\bm\sigma}
    \exp\!\left(t[Z_\ell-\E_{\bm\sigma}Z_\ell]\right)
    \leq
    \exp\!\left(\frac{2t^2B_{\mathfrak C}^2}{\mu}\right).
\]
The log-sum-exp inequality then gives
\[
    \E_{\bm\sigma}\max_{\ell\in[L]}Z_\ell
    \leq
    \max_{\ell\in[L]}\E_{\bm\sigma}Z_\ell
    +\frac{\log L}{t}+\frac{2tB_{\mathfrak C}^2}{\mu}.
\]
Optimizing at
\(t=\sqrt{\mu\log L/(2B_{\mathfrak C}^2)}\) (with the \(L=1\) case
immediate) yields the standard finite-union estimate
\citep{BartlettMendelson2002}
\[
    \hRad_{\cS_\mu}(\cF_{\mathfrak C})
    \leq
    \max_{\ell\in[L]}\hRad_{\cS_\mu}(\cF_\ell)
    +2B_{\mathfrak C}\sqrt{\frac{2\log L}{\mu}}.
\]
The preceding closed-form calculation applies candidate by candidate, so
\[
    \hRad_{\cS_\mu}(\cF_\ell)
    \leq
    \frac{8\Lambda_\ell(\Lambda_\ell+\Upsilon)+2\Upsilon^2}{\sqrt\mu}
    \leq\frac{A_{\mathfrak C}}{\sqrt\mu}.
\]
Substitute these bounds into \cref{eq:finite-union-intermediate} and apply the
candidate-specific inequality
\(\risk_{\infty,\ell}(h)\leq\risk_{w,\ell}(h)+\tau_{w,\ell}\).
This proves \cref{eq:finite-family-bound}. Because the event is simultaneous
over the complete fixed union, it remains valid after any measurable
sample-dependent selector taking values in that union.
\end{proof}

\subsection{Geometrically mixing case}

\begin{proof}
Since \(\mu=N/(2a)\) and \(a\geq1\),
\begin{align*}
    4(\mu-1)\beta_Z(as-w)
    &\leq
    \frac{2N}{a}\beta_0e^{-c(as-w)}\\
    &\leq
    2N\beta_0e^{cw}e^{-cas}.
\end{align*}
The choice in \cref{eq:geometric-block-choice} makes the last expression at
most \(\delta/2\), so \(\delta'\geq\delta/2\). It also gives
\(a=O(\log(N/\delta))\) with the remaining quantities fixed. Since
\(\mu=N/(2a)\), both stochastic terms in
\cref{eq:closed-generalization} are
\(O(\sqrt{a/N})=O(\sqrt{\log(N/\delta)/N})\).

For a nondivisible available sample size, retain
\(N_a=2a\lfloor N/(2a)\rfloor\) observations. When \(N_a\geq2a\), stationarity
permits applying the same bound to this retained consecutive segment, and the
number discarded is at most \(2a-1\). For fixed parameters, the chosen
\(a=O(\log N)\), so \(a\leq N/2\) and \(N_a\geq2a\) for all sufficiently
large \(N\).
\end{proof}

\subsection{Local-Pauli shadow requirement}

\begin{proof}
For independently uniform local \(X/Y/Z\) measurements, a nonidentity
Pauli string of weight \(q\leq k\) has squared shadow norm
\(3^q\leq3^k\). The median-of-means classical-shadow guarantee applied to
\(M\) observables therefore gives
\[
    S
    \geq
    C_{\mathrm{sh}}3^k\varepsilon^{-2}
    \log(2M/\delta_{\mathrm{sh}})
\]
for a universal constant \(C_{\mathrm{sh}}\); see
\citet{HuangKuengPreskill2020}. For \(NR\) reservoir output states, assign failure probability
\(\delta_{\mathrm{sh}}=\delta_{\mathrm{tot}}/(NR)\) to each state and take a
union bound. This gives \cref{eq:shadow-samples-total}. Each state requires
\(S\) independently randomized snapshots, for \(NRS\) state preparations.
Since each snapshot replays the \(w\)-step recursion, the total is \(NRSw\)
channel applications.
\end{proof}

\subsection{Finite-shot feature perturbations}

\begin{proof}
The coordinate-wise event gives
\(\norm{\widehat\Phi_w-\PhiWin}_2\leq\sqrt{R|\cO|}\,\varepsilon\).
The following argument is deterministic and uses only
\(\norm{h}_{\cH_\kappa}\leq\Lambda\), so it holds simultaneously for every
\(h\in\Hball\). The reproducing property and
\cref{lem:matern-stability} yield
\(
\abs{h(\widehat\Phi_w)-h(\PhiWin)}
\leq
\Lambda\norm{\kappa(\widehat\Phi_w,\cdot)-\kappa(\PhiWin,\cdot)}_{\cH_\kappa}
\leq
(\Lambda/\xi)\sqrt{\nu R|\cO|/(\nu-1)}\,\varepsilon
\), proving \cref{eq:shot-prediction}. Both predictions lie in
\([-\Lambda,\Lambda]\); the squared-loss Lipschitz inequality used in the
proof of \cref{cor:loss-truncation} therefore gives \cref{eq:shot-loss}.
\end{proof}
